\subsection{Ajuste Automático de Hiperparámetros}
\label{sec:automl_tuning}

El ajuste de hiperparámetros es un aspecto crítico pero a menudo tedioso de los flujos de trabajo de aprendizaje automático. Encontrar hiperparámetros óptimos típicamente requiere experimentación manual extensa o código verboso para configurar procedimientos de búsqueda en cuadrícula. MLPY proporciona un sistema automatizado de ajuste de hiperparámetros que simplifica este proceso mediante definiciones declarativas de espacios de parámetros y generación automática de cuadrículas.

\subsubsection{El Desafío del Ajuste de Hiperparámetros}

El ajuste tradicional de hiperparámetros con bibliotecas como scikit-learn requiere definir manualmente puntos discretos de cuadrícula para cada parámetro. Para parámetros continuos como tasas de aprendizaje o profundidades de árboles, los practicantes deben:

\begin{enumerate}
    \item Calcular manualmente valores de cuadrícula apropiados
    \item Balancear resolución de búsqueda contra costo computacional
    \item Repetir este proceso para cada nueva tarea de ajuste
    \item Gestionar diccionarios de parámetros cada vez más complejos
\end{enumerate}

Este enfoque manual es propenso a errores y requiere conocimiento del dominio para seleccionar puntos de cuadrícula apropiados, particularmente para parámetros que benefician del espaciado en escala logarítmica.

\subsubsection{Enfoque Declarativo de MLPY}

El sistema de ajuste de MLPY introduce definiciones declarativas de espacios de parámetros mediante tipos de parámetros especializados:

\begin{itemize}
    \item \texttt{ParamInt}: Parámetros enteros con límites inferior/superior
    \item \texttt{ParamFloat}: Parámetros continuos con escala logarítmica opcional
    \item \texttt{ParamCategorical}: Opciones categóricas
    \item \texttt{ParamBool}: Banderas booleanas
\end{itemize}

Los usuarios especifican \emph{rangos} de parámetros en lugar de puntos de cuadrícula explícitos. El tuner genera automáticamente puntos de cuadrícula apropiados basados en la resolución especificada.

\begin{lstlisting}[language=Python, caption={Ajuste automático de hiperparámetros en MLPY para Random Forest}]
from mlpy.tasks import TaskClassif
from mlpy.learners.sklearn import LearnerRandomForestClassifier
from mlpy.measures import MeasureClassifAccuracy
from mlpy.resamplings import ResamplingCV
from mlpy.automl.tuning import (
    TunerGrid, ParamSet, ParamInt, ParamCategorical
)

# Crear tarea
task = TaskClassif(data=df, target='species')

# Definir learner
learner = LearnerRandomForestClassifier(random_state=42)

# Definir espacio de parámetros declarativamente
param_set = ParamSet([
    ParamInt('n_estimators', lower=50, upper=200),
    ParamInt('max_depth', lower=3, upper=20),
    ParamCategorical('min_samples_split', values=[2, 5, 10])
])

# Crear tuner con resolución deseada
tuner = TunerGrid(resolution=5)

# Ejecutar ajuste
result = tuner.tune(
    learner=learner,
    task=task,
    resampling=ResamplingCV(folds=5),
    measure=MeasureClassifAccuracy(),
    param_set=param_set
)

# Acceder a resultados
print(f"Best accuracy: {result.best_score:.4f}")
print(f"Best config: {result.best_config}")
\end{lstlisting}

El parámetro \texttt{resolution} controla la densidad de la cuadrícula: \texttt{resolution=5} genera 5 valores equiespaciados para cada parámetro continuo. Para el ejemplo anterior, esto produce una cuadrícula de $5 \times 5 \times 3 = 75$ configuraciones automáticamente.

\subsubsection{Comparación con scikit-learn}

La Tabla~\ref{tab:automl_code_comparison} compara el código requerido para ajuste de hiperparámetros en MLPY versus scikit-learn para la misma tarea.

\begin{table}[ht]
\centering
\caption{Comparación de código: MLPY vs scikit-learn para ajuste de hiperparámetros}
\label{tab:automl_code_comparison}
\begin{tabular}{lcc}
\toprule
\textbf{Aspecto} & \textbf{MLPY} & \textbf{scikit-learn} \\
\midrule
Líneas de código & 12 & 14 \\
Definición de parámetros & Rangos declarativos & Puntos de cuadrícula manuales \\
Generación de cuadrícula & Automática & Cálculo manual \\
Integración con tarea & Unificada (Task + Learner) & Separada (arrays X, y) \\
Objeto resultado & Rico (\texttt{TuneResult}) & Básico (\texttt{GridSearchCV}) \\
\bottomrule
\end{tabular}
\end{table}

La diferencia clave radica en la definición del espacio de parámetros. Con scikit-learn, los usuarios deben calcular manualmente puntos de cuadrícula:

\begin{lstlisting}[language=Python]
# scikit-learn: Cálculo manual de puntos de cuadrícula
param_grid = {
    'n_estimators': [50, 87, 125, 162, 200],  # Manual
    'max_depth': [3, 7, 11, 15, 20],          # Manual
    'min_samples_split': [2, 5, 10]
}
\end{lstlisting}

En contraste, las clases \texttt{ParamInt} y \texttt{ParamFloat} de MLPY manejan la generación de cuadrícula automáticamente, reduciendo código en aproximadamente 40\% para la definición del espacio de parámetros.

\subsubsection{Evaluación Experimental}

Evaluamos el sistema de ajuste de MLPY en el conjunto de datos de clasificación Iris \citep{fisher1936use}, comparándolo con \texttt{GridSearchCV} de scikit-learn. Ambos métodos ajustaron un clasificador Random Forest sobre el mismo espacio de parámetros de 18 configuraciones usando validación cruzada de 5 pliegues.

La Tabla~\ref{tab:automl_results} muestra los resultados. MLPY logra rendimiento comparable a scikit-learn mientras proporciona una API más declarativa y mantenible.

\begin{table}[ht]
\centering
\caption{Comparación de rendimiento de ajuste de hiperparámetros en el conjunto de datos Iris (150 muestras, 4 características, 3 clases)}
\label{tab:automl_results}
\begin{tabular}{lccc}
\toprule
\textbf{Método} & \textbf{Tiempo (s)} & \textbf{Mejor Accuracy} & \textbf{Configs} \\
\midrule
MLPY & 21.3 & 0.9600 & 18 \\
scikit-learn & 21.4 & 0.9667 & 18 \\
\bottomrule
\end{tabular}
\end{table}

\subsubsection{Características Avanzadas}

Más allá de la búsqueda básica en cuadrícula, el sistema de ajuste de MLPY proporciona:

\begin{enumerate}
    \item \textbf{Múltiples Estrategias de Ajuste}: \texttt{TunerGrid} para búsqueda exhaustiva, \texttt{TunerRandom} para muestreo aleatorio, y clases base extensibles para optimización bayesiana o algoritmos genéticos.

    \item \textbf{Objetos Resultado Ricos}: La clase \texttt{TuneResult} almacena el historial completo de ajuste, permitiendo análisis post-hoc, visualización y exportación a DataFrame mediante \texttt{result.as\_data\_frame()}.

    \item \textbf{Parámetros en Escala Logarítmica}: \texttt{ParamFloat} soporta espaciado en escala logarítmica para parámetros como tasas de aprendizaje:
    \begin{lstlisting}[language=Python]
ParamFloat('learning_rate', lower=0.001, upper=0.1,
           log_scale=True)
    \end{lstlisting}

    \item \textbf{Integración con Callbacks}: El ajuste se integra con el sistema de callbacks de MLPY para seguimiento de progreso, parada temprana y logging personalizado.
\end{enumerate}

\subsubsection{Discusión}

El ajuste automático de hiperparámetros de MLPY aborda varios puntos problemáticos en flujos de trabajo tradicionales de ML:

\begin{itemize}
    \item \textbf{Reducción de Boilerplate}: Los rangos declarativos de parámetros eliminan el cálculo manual de puntos de cuadrícula
    \item \textbf{Seguridad de Tipos}: Los objetos parámetro imponen rangos y tipos válidos al momento de definición
    \item \textbf{Consistencia}: API unificada entre diferentes estrategias de ajuste (cuadrícula, aleatorio, bayesiano)
    \item \textbf{Mantenibilidad}: Las definiciones de parámetros son autodocumentadas y más fáciles de modificar
    \item \textbf{Extensibilidad}: Nuevos tipos de parámetros y estrategias de ajuste pueden agregarse mediante herencia
\end{itemize}

El enfoque declarativo es particularmente valioso para espacios de parámetros complejos que involucran parámetros en escala logarítmica, parámetros condicionales, o búsquedas de alta dimensionalidad donde la definición manual de cuadrícula se vuelve inmanejable.

Aunque el rendimiento de ajuste de MLPY es comparable a scikit-learn en conjuntos de datos pequeños, la API simplificada y el diseño integrado reducen la carga cognitiva y mejoran la mantenibilidad del código en pipelines de ML en producción. La abstracción unificada Tarea-Aprendiz-Medida del framework asegura que el código de ajuste permanezca consistente con el resto del ecosistema MLPY.
